Current practice treats the reliability of a language-model agent as a property of the model and works on it through prompts, instructions, and fine-tuning. This paper treats reliability as a property of the system and takes as its object of design the set of states the agent can reach. A cage is a predicate on transitions that fixes which edges exist; an adjudicator outside the agent's writable space decides whether a reached state is a goal or an unsafe state. A cage is good when the goal region stays reachable and the declared unsafe region does not; the slack between the two conditions is what search consumes. The conditions are independent and behave oppositely under tightening, and six elementary facts about this, including a three-state counterexample in which conjoining two goal-preserving cages loses the goal, are checked in Lean 4 with no external library. From the ontology follow three structural requirements on an adjudicator (non-persuadability, priority, traceability), an eight-step method, and two human positions that cannot be delegated from inside the system. Four hypotheses with numerical decision rules were frozen before any agent run; the paper reports what is already established and marks every outstanding result as pending. The framework generalises supervisory control to a cognitive architecture; its safety statements are relative to a declared model of the unsafe region.
